\scanpage{6}
\mspage{5}

\begin{center}- 56 -\end{center}
\margin{$\mathrm{n}^{\mathrm{o}}357$}

\[
0<\frac{1}{(f_0:f)}\leq\frac{(f:g)}{(f_0:g)}\leq(f:f_0).
\]
\[
\text{(vii)}\qquad(f+f':g)\leq(f:g)+(f':g).
\]
(viii) si $f$, $f'$, $h$ sont dans $\mathcal{K}_+$, si $h\geq1$ sur le support de $f+f'$, et si $\varepsilon>0$, il y a un voisinage compact $V$ de $e$ tel que l'on ait, pour tout $g\in\mathcal{K}^*_+(V)$ :
\[
(f+f':g)\geq(f:g)+(f':g)-\varepsilon(h:g).
\]
Les trois premières sont évidentes ; (iv) aussi, puisque $\sup(\gamma(\mu)g)\leq\Vert\mu\Vert\sup(g)$ ; (v) résulte aussitôt de la formule $\Vert\mu*\mu'\Vert={} =\Vert\mu\Vert.\Vert\mu'\Vert$ (par la vertu du Chap.VII, § 1, n°1, ou, moins savamment et plus simplement dans le cas qui nous occupe ici, par un calcul direct trivial) et de ce qu'on a, si $f\leq\gamma(\mu)g$ et $g\leq\gamma(\mu')h$ :
\[
f\leq\gamma(\mu)\gamma(\mu')h=\gamma(\mu*\mu')h.
\]
Si on applique (v) à $f_0$, $f$, $g$ d'une part, et à $f$, $f_0$, $g$ d'autre part, on obtient (vi) ; (vii) est évidente. Enfin, soient $f$, $f'$, $h$ comme il est dit dans (viii). Posons $F=f+f'+\delta h$, avec $\delta>0$ ; les fonctions $\varphi$, $\varphi'$ qui coïncident respectivement avec $f/F$ et $f'/F$ sur le support de $f+f'$, et qui sont $0$ en dehors de celui-ci, sont dans $\mathcal{K}$, donc uniformément continues ; pour tout $\eta>0$, on peut donc choisir un voisinage compact $V$ de $e$ tel que $|\varphi(x)-\varphi(s)|$ et $|\varphi'(x)-\varphi'(s)|$ soient $\leq\eta$ pour $s^{-1}x\in V$ ; alors, quelle que soit $g\in\mathcal{K}^*_+(V)$, on aura
\[
\varphi(x)g(s^{-1}x)\leq[\varphi(s)+\eta]g(s^{-1}x)
\]
quels que soient $x$ et $s$, et de même pour $\varphi'$. Supposons maintenant qu'on ait, pour tout $x$ :
\[
F(x)\leq\sum_{i=1}^{n}a_i g(s_i^{-1}x).
\]
Puisque $f=\varphi F$, on en conclut :
